% 02-marco-teorico.tex — Fuente: docs/informe-final/02-marco-teorico.md
% Todas las citas marcadas "[PENDIENTE: verificar cita]" en el borrador
% fueron verificadas contra fuente primaria el 17-08-2026 (ver referencias.bib
% y README.md de esta carpeta) y ya no llevan la marca.

\chapter{Marco teórico}
\label{cap:marco-teorico}

Este capítulo se limita a los fundamentos estrictamente necesarios para
comprender las decisiones documentadas en los capítulos posteriores; no es
un compendio general de teoría de aplicaciones web.

\section{Ingeniería de requisitos: ISO/IEC/IEEE 29148:2018 y el marco de Pohl}

La norma ISO/IEC/IEEE~29148:2018~\cite{ISO29148} define los procesos de
ciclo de vida para la ingeniería de requisitos de sistemas y software,
incluyendo la estructura recomendada de un Software Requirements
Specification (SRS) y el patrón sintáctico de enunciado de requisito
\emph{[condición] [sujeto] shall [acción] [objeto] [restricción]}.
Pohl~\cite{POHL2010} organiza el proceso en cuatro actividades
interdependientes ---elicitación, negociación, documentación y
validación--- que este proyecto usa como estructura del
Capítulo~\ref{cap:ingenieria-requisitos}. Wiegers y
Beatty~\cite{WIEGERS2013} complementan el marco con las prácticas de
desarrollo y gestión de requisitos (priorización, control de cambios)
usadas en \texttt{docs/requisitos/CHANGELOG-REQ.md}.

\section{Calidad de requisitos: INCOSE Guide to Writing Requirements v4}

La guía INCOSE~v4~\cite{INCOSE2023} define nueve características de
calidad para un requisito individual (C1--C9: Necessary, Appropriate,
Unambiguous, Complete, Singular, Feasible, Verifiable, Correct,
Conforming) y seis para un conjunto de requisitos (C10--C15: Complete,
Consistent, Feasible, Comprehensible, Able to be validated, Correct). Este
proyecto aplica ambos conjuntos de características en
\texttt{docs/checklists/incose-checklist.md}, con evidencia citada
requisito por requisito en vez de una declaración genérica de
cumplimiento.

\section{Historias de usuario y casos de uso como técnicas complementarias}

Las historias de usuario en formato Connextra (\emph{As a [rol], I want
[funcionalidad], so that [beneficio]}) con criterios INVEST (Independent,
Negotiable, Valuable, Estimable, Small, Testable) capturan la intención
funcional desde la perspectiva del usuario, mientras que los casos de uso
en la plantilla de Cockburn~\cite{COCKBURN2000} formalizan el flujo de
interacción en niveles 1 a 4 (resumen, cometido de usuario, subfunción,
implementación). El proyecto usa ambas técnicas de forma complementaria,
no sustitutiva: las historias en \texttt{docs/requisitos/historias/}
capturan el valor de negocio; los casos de uso en
\texttt{docs/requisitos/casos-de-uso/} capturan el flujo detallado de
interacción para los escenarios críticos.

\section{Modelo arquitectónico C4}

El modelo C4 de Brown~\cite{BROWN2018} describe la arquitectura de un
sistema en cuatro niveles progresivos de detalle ---Contexto, Contenedores,
Componentes y Código--- cada uno dirigido a una audiencia distinta (desde
stakeholders no técnicos hasta desarrolladores). Este proyecto documenta
los niveles 1 a 3 en \texttt{docs/diagramas/}
(\texttt{C4\_Nivel1\_Contexto.md}, \texttt{C4\_Nivel2\_Contenedores.md},
\texttt{C4\_Nivel3\_Componentes\_Backend.md}), usados como base del
Capítulo~\ref{cap:diseno-arquitectura}.

\section{Calidad de software: ISO/IEC 25010}

ISO/IEC~25010:2011~\cite{ISO25010} define un modelo de características de
calidad de producto software (adecuación funcional, eficiencia de
desempeño, compatibilidad, usabilidad, fiabilidad, seguridad,
mantenibilidad, portabilidad) usado en este proyecto para clasificar los
requisitos no funcionales del SRS y para estructurar la tabla de atributos
de calidad del Capítulo~\ref{cap:diseno-arquitectura}. Las dimensiones
evaluadas empíricamente en el Capítulo~\ref{cap:evaluacion-resultados}
(rendimiento, seguridad, usabilidad) se mapean directamente a
subcaracterísticas de este modelo.

\section{Principios REST}

La arquitectura REST (Representational State Transfer), formalizada por
Fielding en su disertación doctoral~\cite{FIELDING2000}, define un
conjunto de restricciones arquitectónicas (interfaz uniforme, sin estado
en el servidor, cacheable, sistema en capas) que la API de Artisync sigue
para exponer sus recursos (\texttt{/api/v1/...}), documentada con OpenAPI
(Springdoc).

\section{Autenticación: JWT (RFC 7519)}

JSON Web Token (RFC~7519)~\cite{JWT2015} especifica un formato compacto y
auto-contenido para representar afirmaciones (\emph{claims}) entre dos
partes, firmado digitalmente para garantizar integridad. Artisync usa JWT
firmado HS256 para el access token; la decisión de dónde alojarlo (cuerpo
de respuesta vs. cookie \texttt{HttpOnly}) y sus alternativas descartadas
se documentan en \texttt{docs/adr/adr-002-esquema-autenticacion.md}.

\section{Controles OWASP Top 10}

El OWASP Top 10:2021~\cite{OWASP2021} cataloga las diez categorías de
riesgo de seguridad más críticas en aplicaciones web, resultado de un
análisis de datos de la industria. Este proyecto evidencia seis de esas
categorías (A01 Control de acceso roto, A02 Fallas criptográficas, A03
Inyección, A05 Configuración de seguridad incorrecta, A07 Fallas de
identificación y autenticación, A09 Fallas de registro y monitoreo de
seguridad) con evidencia reproducible en \texttt{docs/mediciones/sec/},
complementadas con un escaneo automatizado OWASP ZAP baseline y análisis
estático SpotBugs + find-sec-bugs.

\section{Patrones de acceso a datos y modelos de consistencia}

La disyuntiva entre mapeo objeto-relacional (ORM) y acceso directo
mediante SQL/procedimientos almacenados es un compromiso reconocido en la
literatura de arquitectura de datos y de ingeniería de software empírica
sobre seguridad: el ORM reduce el código repetitivo de las operaciones
CRUD elementales a costa de perder control fino sobre consultas complejas,
mientras que los procedimientos almacenados ofrecen ejecución más cercana
a los datos y una superficie de ataque de inyección SQL distinta
(mitigable con parametrización estricta). La literatura fundacional sobre
detección y prevención de inyección SQL --- tanto por chequeo estático de
consultas generadas dinámicamente~\cite{GOULD2004} como por monitoreo en
tiempo de ejecución de un modelo de consultas legítimas~\cite{HALFOND2005}
--- coincide en que el riesgo real no es el mecanismo elegido (ORM vs.
procedimiento almacenado) sino la concatenación de entrada de usuario en
el texto del comando SQL, sin importar cuál de los dos mecanismos la
introduzca. La guía OWASP de prevención de inyección
SQL~\cite{OWASPSQLI} reconoce ambos mecanismos ---consultas parametrizadas
vía ORM y procedimientos almacenados correctamente parametrizados--- como
defensas primarias equivalentes contra la inyección, siempre que ninguno
de los dos concatene entrada de usuario en el texto del comando SQL. Esta
equivalencia es la base técnica de la estrategia híbrida adoptada en
\texttt{docs/adr/adr-006-estrategia-acceso-datos.md} y evaluada en el
Capítulo~\ref{cap:evaluacion-resultados}.

\section{Documentación de decisiones de arquitectura (ADR)}

El formato de \emph{Architecture Decision Record} propuesto por
Nygard~\cite{NYGARD2011} --- un documento corto por decisión, con
contexto, decisión y consecuencias --- se adoptó en este proyecto
precisamente por la razón que motivó su creación: los documentos grandes
de arquitectura no se mantienen actualizados, mientras que los documentos
pequeños y modulares sí tienen una oportunidad real de actualizarse. Los
siete ADR de Artisync (\texttt{docs/adr/}) siguen esta plantilla; el
Capítulo~\ref{cap:diseno-arquitectura} resume su contenido.
